
\documentclass{article}
\usepackage{colortbl}
\usepackage{makecell}
\usepackage{multirow}
\usepackage{supertabular}

\begin{document}

\newcounter{utterance}

\centering \large Interaction Transcript for game `wordle', experiment `high\_frequency\_words\_no\_clue\_no\_critic', episode 9 with reflexion{-}llama{-}3.1{-}8b{-}instruct.
\vspace{24pt}

{ \footnotesize  \setcounter{utterance}{1}
\setlength{\tabcolsep}{0pt}
\begin{supertabular}{c@{$\;$}|p{.15\linewidth}@{}p{.15\linewidth}p{.15\linewidth}p{.15\linewidth}p{.15\linewidth}p{.15\linewidth}}
    \# & \multicolumn{2}{c}{Player} && \multicolumn{2}{c}{Game Master} \\
    \hline

    \theutterance \stepcounter{utterance}  
    & & & \multicolumn{4}{p{0.6\linewidth}}{
        \cellcolor[rgb]{0.9,0.9,0.9}{
            \makecell[{{p{\linewidth}}}]{
                \texttt{\tiny{[P1$\langle$GM]}}
                \texttt{You have attempted this same task before and did not succeed. These are the plans you wrote after those attempts.} \\
\texttt{Trial \#1: New plan: } \\
\\ 
\texttt{In my last attempt, I made a mistake by producing a guess that was not a valid five{-}letter word. I also failed to follow the format by not including the "explanation:" tag. Next time, I will make sure to produce a valid five{-}letter word that includes the letters that were previously yellow (l and e) and excludes the letters that were previously red (h, o, u, s, g, o, b). I will also follow the format by including the "explanation:" tag and making sure that my guess is a valid five{-}letter word. For example, I could try a word like "leapt".} \\
\\ 
\texttt{Trial \#2: New plan: } \\
\\ 
\texttt{In my last attempt, I made a mistake by not paying attention to the fact that the letter "e" was already in the correct position in the word. I kept trying to find a new position for the letter "e" instead of focusing on the other letters that were still yellow (l) or red (k, s, c, h, g, a, m, f, p, t, o, u, and b). Next time, I will focus on finding a new word that includes the letters "l" and "e" in the correct position and excludes the letters that were previously red. I will also try to find a word that includes the letter "e" in a different position, but not necessarily in the last position. For example, I could try a word like "leaks" or "leens".} \\
\\ 
\texttt{Trial \#3: New plan: } \\
\\ 
\texttt{In my last attempt, I made a mistake by producing a guess that was not a valid five{-}letter word. I also failed to follow the format by not including the "explanation:" tag. Next time, I will make sure to produce a valid five{-}letter word that includes the letters that were previously yellow (l and e) and excludes the letters that were previously red (a, p, t). I will also follow the format by including the "explanation:" tag and making sure that my guess is a valid five{-}letter word. For example, I could try a word like "lemon" or "leech".} \\
\\ 
\\ 
\texttt{You are a language wizard who likes to guess words by using the given rules.} \\
\\ 
\texttt{Welcome to Wordle! You have six attempts to guess the target word, a valid English word of five lowercase letters (a{-}z). Please use the tags "explanation:" and "guess:" to provide a concise explanation for each guess.} \\
\\ 
\texttt{For instance, if your guess is "apple", your response should be} \\
\texttt{explanation: this is a common five{-}letter English word, and I am starting my guess with this word.} \\
\texttt{guess: apple} \\
\\ 
\texttt{After each guess, your answer will be validated, and you will receive feedback indicating which letters are correct (green), which letters are correct but in the wrong position (yellow), and which letters are incorrect (red). This feedback can be useful in determining which letters to include or exclude in your next guess.} \\
\\ 
\texttt{For example, the feedback for "apple" might be:} \\
\texttt{guess\_feedback: a<yellow> p<yellow> p<green> l<yellow> e<red>} \\
\\ 
\texttt{The explanation should contain details about how the guess\_feedback is used to arrive at a new guess.} \\
\\ 
\texttt{Do not generate any other text.} \\
\\ 
\texttt{Let's begin with your first guess.} \\
            }
        }
    }
     \\ \\

    \theutterance \stepcounter{utterance}  
    & \multicolumn{4}{p{0.6\linewidth}}{
        \cellcolor[rgb]{0.561,0.737,0.561}{
            \makecell[{{p{\linewidth}}}]{
                \texttt{\tiny{[P1$\rangle$GM]}}
                \texttt{explanation: I will start by guessing a common five{-}letter English word that includes the letters that were previously yellow (l and e) and excludes the letters that were previously red (h, o, u, s, g, b).} \\
\texttt{guess: leapt} \\
            }
        }
    }
    & & \\ \\

    \theutterance \stepcounter{utterance}  
    & & & \multicolumn{2}{p{0.3\linewidth}}{
        \cellcolor[rgb]{0.9,0.9,0.9}{
            \makecell[{{p{\linewidth}}}]{
                \texttt{\tiny{[GM$|$GM]}}
                \texttt{attempts = 1} \\
\texttt{target\_word = level} \\
\texttt{guess = leapt} \\
\texttt{guess\_feedback = l<green> e<green> a<red> p<red> t<red>} \\
            }
        }
    }
    & & \\ \\

    \theutterance \stepcounter{utterance}  
    & & & \multicolumn{4}{p{0.6\linewidth}}{
        \cellcolor[rgb]{0.9,0.9,0.9}{
            \makecell[{{p{\linewidth}}}]{
                \texttt{\tiny{[P1$\langle$GM]}}
                \texttt{guess\_feedback: l<green> e<green> a<red> p<red> t<red>} \\
\\ 
\texttt{Provide your response only in this format.} \\
\texttt{explanation: details} \\
\texttt{guess: word} \\
            }
        }
    }
     \\ \\

    \theutterance \stepcounter{utterance}  
    & \multicolumn{4}{p{0.6\linewidth}}{
        \cellcolor[rgb]{0.561,0.737,0.561}{
            \makecell[{{p{\linewidth}}}]{
                \texttt{\tiny{[P1$\rangle$GM]}}
                \texttt{explanation: Since the letters "l" and "e" are both green, I will keep them in the correct position. The letters "a", "p", and "t" are all red, so I will exclude them from my next guess. I will also try to find a word that includes the letter "l" and "e" in the correct position and excludes the letters that were previously red.} \\
\texttt{guess: length} \\
            }
        }
    }
    & & \\ \\

    \theutterance \stepcounter{utterance}  
    & & & \multicolumn{2}{p{0.3\linewidth}}{
        \cellcolor[rgb]{0.9,0.9,0.9}{
            \makecell[{{p{\linewidth}}}]{
                \texttt{\tiny{[GM$|$GM]}}
                \texttt{Error: The length of the guessed word is not 5.} \\
            }
        }
    }
    & & \\ \\

    \theutterance \stepcounter{utterance}  
    & & & \multicolumn{2}{p{0.3\linewidth}}{
        \cellcolor[rgb]{0.9,0.9,0.9}{
            \makecell[{{p{\linewidth}}}]{
                \texttt{\tiny{[GM$|$GM]}}
                \texttt{game\_result = ABORT} \\
            }
        }
    }
    & & \\ \\

    \theutterance \stepcounter{utterance}  
    & & \multicolumn{4}{p{0.6\linewidth}}{
        \cellcolor[rgb]{0.95,0.95,0.95}{
            \makecell[{{p{\linewidth}}}]{
                \texttt{\tiny{[P1$\leftrightarrow$P]}}
                \texttt{Post{-}episode Reflexion (not sent to GM):} \\
\texttt{In my last attempt, I made a mistake by producing a guess that was not a valid five{-}letter word. I also failed to follow the format by not including the "explanation:" tag. Next time, I will make sure to produce a valid five{-}letter word that includes the letters that were previously yellow (l and e) and excludes the letters that were previously red (a, p, t). I will also follow the format by including the "explanation:" tag and making sure that my guess is a valid five{-}letter word. Additionally, I will pay closer attention to the fact that the letter "e" is already in the correct position in the word and focus on finding a new word that includes the letters "l" and "e" in the correct position and excludes the letters that were previously red. For example, I could try a word like "lemon" or "leech".} \\
            }
        }
    }
    & \\ \\

\end{supertabular}
}

\end{document}
